%% main.tex --- KIISE 논문지 양식 사용 예시
%% 컴파일: xelatex main.tex
%%
%% 심사용:  \documentclass{kiise}        (저자 정보 일체 미표기, 감사의 글 미기재)
%% 채택 후: \documentclass[final]{kiise} (저자, 약력 출력)
\documentclass[final]{kiise}

%%%%%%%%%%%%%%%%%%%%%%%%%%%%%%%%%%%%%%%%%%%%%%%
% include additional packages you need to use
%%%%%%%%%%%%%%%%%%%%%%%%%%%%%%%%%%%%%%%%%%%%%%%
\usepackage{appendix}

% graphic, float package
\usepackage{graphicx}		% for setting images
\usepackage{float}			% for float objects
% \usepackage{subfloat}		% subfloat.sty 미설치 → 비활성화 (미사용)
\usepackage{subfigure}		% for adding several figures in a figure environment
\usepackage{lscape}			% for landscape type images or tables

\usepackage{enumitem}
\usepackage{xspace}			% for \model, \eg 등 매크로 trailing space 처리
\usepackage[table]{xcolor}	% tikz보다 먼저 [table] 옵션으로 로드 (옵션 충돌 방지)
\usepackage{tikz}			% 모델 구조 그림
\usetikzlibrary{arrows.meta, positioning, fit, backgrounds}

% for compact section title spacing
% \usepackage[compact]{titlesec}


% mathmetical presentation
\usepackage{gensymb}
\usepackage{amsmath}
\usepackage{amssymb}
\usepackage{amsthm}
\usepackage{exscale}
\usepackage{textcomp}		% extra symbols


% for circled number
\newcommand{\cl}[1]{\textcircled{\scriptsize #1}}


% package for using algorithmic presentation (algorithmic.sty 미설치 → 비활성화)
% \usepackage{algorithmic}
% \usepackage{algorithm}
% customize algorithmic environment
% \renewcommand{\algorithmicrequire}{\makebox[40px]{\hfill\textbf{Input :}}}
% \renewcommand{\algorithmicensure}{\makebox[40px]{\hfill\textbf{Output :}}}

% array and table presentation
\usepackage{array}
\usepackage{tabulary}
\usepackage{multirow}
% \usepackage[table]{xcolor}  % 위에서 tikz 앞에 이미 로드함 (중복 제거)
\usepackage{ctable}
\usepackage{booktabs}		% for typesetting tables at the level of publication
\usepackage{refcount}		% \getrefnumber: 표/그림 번호 발음 기반 자동 조사용     % 추가 패키지
% 공통 매크로 정의
\newcommand{\model}{CSDG\xspace}
\newcommand{\modelfull}{Cluster-Sequence Generation\xspace}
\newcommand{\eg}{e.g.,\xspace}
\newcommand{\ie}{i.e.,\xspace}

% ===== 표/그림 번호 발음 기반 자동 조사 =====
% \josaref{표}{label}{UN}  ->  [표 N] + 은/는 (번호가 바뀌어도 자동 일치)
%   UN: 은/는, I: 이/가, UL: 을/를, GW: 과/와
% 끝자리 발음에 받침이 있으면(0,1,3,6,7,8) 받침용 조사, 없으면(2,4,5,9) 모음용 조사.
\makeatletter
\newif\if@jong
\def\jUN{UN}\def\jI{I}\def\jUL{UL}\def\jGW{GW}
\newcommand{\josaref}[3]{%
  [#1~\ref{#2}]%
  \@tempcnta=\getrefnumber{#2}\relax
  \@tempcntb=\@tempcnta \divide\@tempcntb by 10 \multiply\@tempcntb by 10
  \advance\@tempcnta by -\@tempcntb % 끝자리
  \@jongfalse
  \ifnum\@tempcnta=0 \@jongtrue\fi
  \ifnum\@tempcnta=1 \@jongtrue\fi
  \ifnum\@tempcnta=3 \@jongtrue\fi
  \ifnum\@tempcnta=6 \@jongtrue\fi
  \ifnum\@tempcnta=7 \@jongtrue\fi
  \ifnum\@tempcnta=8 \@jongtrue\fi
  \def\@jt{#3}%
  \if@jong
    \ifx\@jt\jUN 은\else\ifx\@jt\jI 이\else\ifx\@jt\jUL 을\else\ifx\@jt\jGW 과\fi\fi\fi\fi
  \else
    \ifx\@jt\jUN 는\else\ifx\@jt\jI 가\else\ifx\@jt\jUL 를\else\ifx\@jt\jGW 와\fi\fi\fi\fi
  \fi
}
\makeatother
  % 공통 매크로 (\model, \eg, \josaref 등)

\begin{document}

%% ===== 논문 구성 순서 (규정 7-(1)) =====
%% 국문 제목 -> 영문 제목 -> 국문 요약 -> 국문 키워드 -> 영문 요약 -> 영문 키워드
%% -> 본문 -> 참고문헌 -> 부록

\koreantitle{정보과학회 저널 LaTeX 양식 만들었습니다\\ 제목을 입력하세요}
\englishtitle{NLP Lab @ Hanbat National University}

%% 채택 후(final)에만 출력됨 — 심사용에서는 무시
%% 머리글: 홀수면=짧은 제목, 짝수면=학회지 권호 정보
\runningtitle{국립한밭대학교 NLP Lab }
\journalinfo{정보과학회논문지: 정보통신 제 31 권 제 6 호(2004.12)}

\koreanauthors{홍 길 동 \qquad 김 철 수}
\englishauthors{Gildong Hong \qquad Cheolsu Kim}

%% 1면 좌측단 각주: 사사, 저자 소속, 교신저자, 접수/심사일 (final 모드에서만 출력)
\firstpagefootnotes{%
  \fpitem{\dag}{
  사사 입력
  }
  \fpitem{\ddag}{정 회 원 : 국립한밭대학교 컴퓨터공학과\\ \mbox{}\quad parkce@hanbat.ac.kr}
  \fpitem{\S}{정 회 원 :  \\
  \mbox{}\quad 교신저자\\ \mbox{}\quad parkce@hanbat.ac.kr}
  \fpitem{}{논문접수 : }
  \fpitem{}{심사완료 : }
}

%% 국문 요약: 300~500자 (규정 7-(2))
\koreanabstract{요약을 입력하세요. 300--500자.}

%% 국문 키워드: 4~6개
\koreankeywords{정보과학회 저널, Latex 양식}

%% 영문 요약: 100~200 단어
\englishabstract{
English Abstract
}

%% 영문 키워드: 4~6개
\englishkeywords{Proxy signature, ID-based partially blind signature, Proxy
partially blind signature, Gap Diffie-Hellman problem, Bilinear-pairing}

\maketitle



\section{서 론}

대리서명 방식은 원서명자를 대신해서 서명을 할 수 있는 대리서명자를 지정하여
대신 서명하도록 하는 서명 방식이다. 본문에 직접 관련이 있는 문헌은 본문 중에
참고문헌 번호를 쓰고 인용 순서대로 참고문헌란에 기술한다\cite{kim2008,nguyen2008}.


\section{관련 연구}
절 번호는 1.1과 같이 아라비아 숫자로 표기한다(규정 9). 수식은 다음과 같이
번호를 붙여 작성한다.
\begin{equation}
  e(aP, bQ) = e(P, Q)^{ab}
\end{equation}

\section{본 론}
~
%% ----- 표 예시: 제목은 상단 중앙 -----
\begin{table}[t]
\tabcaption{제안 방식의 계산량 비교}{Comparison of computational costs}\label{tab:cost}
\centering
\fontsize{8pt}{11pt}\selectfont
\begin{tabular}{lccc}
\toprule
Scheme & Pairing & Scalar Mul. & Hash \\
\midrule
Zhang et al. & 3 & 4 & 2 \\
Proposed     & 2 & 3 & 2 \\
\bottomrule
\end{tabular}
\end{table}

본 논문에서 제안하는 \model(\modelfull) 기법을 설명한다.
\josaref{표}{tab:cost}{UN} 계산량 비교를 보여주며 (\eg pairing 횟수),
그림과 표는 인쇄될 크기의 두 배 정도로 선명하게 작성하여 본문 중에 위치를
정한다. 그림 및 표의 내용은 모두 영문으로 작성하고, 제목은 한글과 영문을
병기한다(규정 10).


%% ----- 그림 예시: 제목은 하단 중앙 -----
\begin{figure}[t]
\centering
\fbox{\rule{0pt}{2.5cm}\rule{0.85\columnwidth}{0pt}}% 실제 그림으로 교체
\figcaption{제안하는 서명 프로토콜의 전체 구조}{Overall architecture of the proposed signature protocol}
\end{figure}

\section{실 험}

\section{결론}
단편논문은 출판 논문지의 편집형식 기준 5면 이내로 작성한다(규정 7-(3)).
심사용 논문에는 감사의 글을 기재하지 않는다(규정 7-(4)).


%% ===== 참고문헌: BibTeX 사용 (영문 표기, 인용 순서, 규정 8) =====
\bibliographystyle{kiise}
\bibliography{references}

%% ===== 저자 약력 (채택 후, final 옵션에서만 출력; 200자 이내, 규정 11) =====
\biography{홍 길 동}{}{2010년 한밭대학교 컴퓨터공학과 졸업(학사). 2012년 동
대학원 졸업(석사). 2012년$\sim$현재 동 대학원 박사과정. 관심분야는 암호이론,
정보보호}

\end{document}
